\clearpage\thispagestyle{empty}
\pdfbookmark[0]{Sources for Bright Spot on the Moon}{moonmarch1821-sources}
\begin{center}\textbf{Sources for Bright Spot on the Moon}\end{center}
\noindent This AI translation follows the report [\textit{Heller Fleck im Mond.}] in \href{https://gdz.sub.uni-goettingen.de/id/PPN236007467}{\textit{Werke}, vol.~VI} (Göttingen, 1874), pp.~436--437. The report quotes Dr. Heinrich Wilhelm Olbers’s letter to Gauss dated 27 February 1821, including the diary entry, and refers to a later letter from Olbers. The first printing appeared under Göttingen in \textit{Göttingische gelehrte Anzeigen}, no.~46, 22 March 1821, pp.~\href{https://images.p.sub.uni-goettingen.de/iiif/image/gdz:PPN319721507_1821:00000451/full/1600,/0/default.jpg}{449}, \href{https://images.p.sub.uni-goettingen.de/iiif/image/gdz:PPN319721507_1821:00000452/full/1600,/0/default.jpg}{450}, \href{https://images.p.sub.uni-goettingen.de/iiif/image/gdz:PPN319721507_1821:00000453/full/1600,/0/default.jpg}{451} and \href{https://images.p.sub.uni-goettingen.de/iiif/image/gdz:PPN319721507_1821:00000454/full/1600,/0/default.jpg}{452}. The digitized copy is from the Göttingen State and University Library.
\par\smallskip\noindent In the comparison with a fixed star northeast of the Moon, the collected text gives third magnitude and the first printing on p.~\href{https://images.p.sub.uni-goettingen.de/iiif/image/gdz:PPN319721507_1821:00000452/full/1600,/0/default.jpg}{450} gives sixth. The later hypothetical brightness of sixth magnitude occurs in both witnesses. The collected \textit{Maeilius} differs from the apparent reading \textit{Manilius} on p.~\href{https://images.p.sub.uni-goettingen.de/iiif/image/gdz:PPN319721507_1821:00000454/full/1600,/0/default.jpg}{452} of the first printing. The translation follows the collected text. Translation and typesetting produced by AI.
\clearpage
